% Three scatter-plot panels distinguishing independence, zero
% correlation but dependent, and strong positive correlation.
\begin{figure}[htbp]
\centering
\begin{tikzpicture}[x=1cm, y=1cm, font=\small]

% ---- Panel 1: Independent, both standard normal ----
\begin{scope}[shift={(0, 0)}]
  \node[anchor=south, font=\small\bfseries] at (2.0, 3.3)
    {Independent};
  \node[anchor=south, font=\scriptsize] at (2.0, 2.9)
    {$\rho = 0$, $X \perp Y$};
  \draw[->, thick] (-0.1, 0) -- (4.3, 0);
  \draw[->, thick] (0, -0.1) -- (0, 2.7);
  \draw[thin, enggray] (2.0, 0) -- (2.0, 2.5);
  \draw[thin, enggray] (0, 1.25) -- (4.0, 1.25);
  % Cloud of stylised points
  \foreach \x/\y in {1.4/1.0, 2.5/1.5, 1.8/1.6, 2.3/0.9, 2.7/1.4,
                     1.6/1.3, 2.4/1.8, 2.1/0.7, 1.9/1.9, 2.5/1.1,
                     1.5/1.5, 2.6/1.7, 1.7/0.8, 2.2/1.4, 2.8/1.3,
                     1.4/1.2, 2.3/1.6, 2.0/1.0, 1.8/1.0, 2.4/1.4} {
    \fill[engblue] (\x, \y) circle (0.045);
  }
  \node[anchor=north, font=\scriptsize] at (2.0, -0.1) {$X$};
\end{scope}

% ---- Panel 2: Zero correlation but dependent (Y = X^2) ----
\begin{scope}[shift={(5.2, 0)}]
  \node[anchor=south, font=\small\bfseries] at (2.0, 3.3)
    {Dependent, $\rho = 0$};
  \node[anchor=south, font=\scriptsize] at (2.0, 2.9)
    {$Y = X^{2}$ on $X \sim \mathcal{N}(0,1)$};
  \draw[->, thick] (-0.1, 0) -- (4.3, 0);
  \draw[->, thick] (0, -0.1) -- (0, 2.7);
  % U-shape on standardised axes; X in [-2.5, 2.5] mapped to [0, 4]
  \foreach \x in {-2.4, -2.0, -1.6, -1.2, -0.8, -0.4, 0.0,
                  0.4, 0.8, 1.2, 1.6, 2.0, 2.4} {
    \pgfmathsetmacro\px{2.0 + 0.8*\x}
    \pgfmathsetmacro\py{0.25*\x*\x}
    \fill[engred] (\px, \py) circle (0.045);
  }
  \node[anchor=north, font=\scriptsize] at (2.0, -0.1) {$X$};
\end{scope}

% ---- Panel 3: Strong positive correlation ----
\begin{scope}[shift={(10.4, 0)}]
  \node[anchor=south, font=\small\bfseries] at (2.0, 3.3)
    {Correlated, $\rho = 0.85$};
  \node[anchor=south, font=\scriptsize] at (2.0, 2.9)
    {bivariate Gaussian};
  \draw[->, thick] (-0.1, 0) -- (4.3, 0);
  \draw[->, thick] (0, -0.1) -- (0, 2.7);
  % Stylised positive-slope cloud
  \foreach \x/\y in {0.8/0.6, 1.0/0.9, 1.2/1.0, 1.4/1.1, 1.5/1.3,
                     1.7/1.4, 1.8/1.5, 2.0/1.6, 2.1/1.4, 2.2/1.7,
                     2.4/1.8, 2.5/2.0, 2.7/2.1, 2.9/2.3, 3.1/2.2,
                     3.2/2.4, 1.3/0.7, 1.6/1.0, 2.6/2.1, 0.9/0.7} {
    \fill[engamber] (\x, \y) circle (0.045);
  }
  % Regression line, slope ~0.85
  \draw[thick, engblue, dashed] (0.4, 0.2) -- (3.5, 2.5);
  \node[anchor=north, font=\scriptsize] at (2.0, -0.1) {$X$};
\end{scope}

\end{tikzpicture}
\caption[Independence, zero correlation, and linear dependence]{Three
distinct joint structures with informative Pearson correlations.
Left: independent $X$ and $Y$, each standard normal; the cloud is
isotropic. Centre: $Y = X^{2}$ with $X \sim \mathcal{N}(0, 1)$; the
correlation is zero by symmetry, yet $Y$ is a deterministic
function of $X$. Right: bivariate Gaussian with $\rho = 0.85$; the
linear trend is visible and the dashed regression line summarises
it. Zero correlation does not imply independence; the converse
holds only for the joint Gaussian. The diagnostic implication for
working with data: a Pearson correlation near zero rules out a
linear relationship and nothing more. Source: computed by the
authors; the joint Gaussian-uncorrelated equivalence is
\cite[ch.~4]{text:wasserman2004}.}
\label{fig:vol02:ch13:independence-correlation}
\end{figure}
